\section{模型建立与求解}

\subsection{电量平衡与储能运行的基本模型}

四问均需保证每一时段供电充足，且电池在连续充放电后仍满足功率与储电量限制。因此，先建立共同的物理模型，再分别规定购电计划的形成时刻、调整方式和费用。

\subsubsection{电量平衡与设备约束}

原始附件以功率记录负荷和光伏出力。乘以时段长度后，得到相应电量及二者之差：
\begin{equation}
  \ell_t=P_t^{\rm load}\Delta t,\qquad
  v_t=P_t^{\rm pv}\Delta t,\qquad n_t=\ell_t-v_t.
  \label{eq:net-load}
\end{equation}
其中，$n_t$称为净负荷电量。$n_t>0$表示光伏不足以满足负荷，$n_t<0$表示光伏有剩余。模型同时保留$\ell_t$与$v_t$，以分别计算光伏利用量和弃电量。

计划购电量决定需要支付的费用，实际取用量还取决于运行时的供需。设当前有效计划为$q_t$，其中实际取用$h_t$、未利用$w_t$，光伏弃电量为$z_t$，紧急购电量为$e_t$，则有
\begin{equation}
  0\le h_t\le q_t,\qquad w_t=q_t-h_t,\qquad
  0\le z_t\le v_t,\qquad e_t\ge0.
  \label{eq:energy-bounds}
\end{equation}
外网取电、光伏和电池放电共同承担负荷及电池充电，扣除弃电后，时段电量平衡为
\begin{equation}
  h_t+e_t+v_t+d_t=\ell_t+c_t+z_t.
  \label{eq:energy-balance}
\end{equation}
式中$e_t$表示实际购入的补充电量，需按紧急电价付费。设置该变量后，供能不足由外网补足，不能将它仅作为允许缺供的罚项。

充入电池的$c_t$只有$\eta_c c_t$转化为内部储电量；向外输出$d_t$则需消耗$d_t/\eta_d$的内部电量。因此，储电量按下式变化：
\begin{equation}
  E_{t+1}=E_t+\eta_cc_t-\frac{d_t}{\eta_d}.
  \label{eq:soc-transition}
\end{equation}
题给额定容量为$12000\ \mathrm{kWh}$。考虑避免过充过放的要求，各时段起止电量必须满足
\begin{equation}
  E^{\min}=1200\le E_t\le10800=E^{\max},
  \qquad t=0,\ldots,T.
  \label{eq:soc-bounds}
\end{equation}
由此可用的储电量范围宽度为$9600\ \mathrm{kWh}$。最大充、放电功率均为$5000\ \mathrm{kW}$，换算成每10分钟的电量上限为
\begin{equation}
  \bar c=\bar d=5000\Delta t=\frac{2500}{3}\ \mathrm{kWh}.
  \label{eq:power-limit}
\end{equation}
用$u_t\in\{0,1\}$表示充放电状态。$u_t=1$时允许充电，$u_t=0$时允许放电，从而得到互斥约束
\begin{equation}
  0\le c_t\le\bar c u_t,\qquad
  0\le d_t\le\bar d(1-u_t).
  \label{eq:exclusive}
\end{equation}

按照紧急购电只用于补足当前负荷的约定，问题二至四还需满足
\begin{equation}
  e_tc_t=e_tz_t=0,\qquad
  e_t=\bigl[\ell_t+c_t+z_t-v_t-d_t-q_t\bigr]^+,
  \label{eq:emergency-operation}
\end{equation}
其中$[x]^+=\max(x,0)$。具体执行时，若供能不足，先停止充电，再购买缺少的电量，避免一边高价补购、一边充电或弃光。将式~\eqref{eq:energy-bounds}--\eqref{eq:emergency-operation}所规定的全部运行条件记为可行域$\mathcal F$。

将各时段电量平衡及储电量变化相加，可得总实际购电量与系统电量的关系：
\begin{equation}
\begin{aligned}
  \sum_{t=0}^{T-1}(h_t+e_t)
  ={}&\sum_{t=0}^{T-1}(\ell_t-v_t+z_t)+(E_T-E_0)\\
     &+\sum_{t=0}^{T-1}\left[(1-\eta_c)c_t+
       \left(\frac{1}{\eta_d}-1\right)d_t\right].
\end{aligned}
\label{eq:aggregate-balance}
\end{equation}
该式表明，实际购电量除补足净负荷外，还需承担弃光造成的缺口、充放电损耗及期末储电量的增加。因此，储能虽然能改变购电时段，也会改变总购电量，费用比较必须同时计算损耗和剩余电量。

\subsubsection{可用信息与跨日储电量}

设$\mathcal I_t$为采取动作时已经获得的信息，包括历史实测、已发布预报、可见价格、当前储电量及有效购电计划。用$\pi_t$表示依据这些信息选择动作的规则，则
\begin{equation}
  a_t=\pi_t(\mathcal I_t).
  \label{eq:causality}
\end{equation}
这要求每次规划只使用当时已经获得的数据。储能在时段内的响应则依据已测供需，具体离散近似在问题二中说明。

例如，每天0时只能利用此前历史和当时已发布的预报制定计划，不能把当天完整实测曲线作为预测输入。若两种未来情景在当前具有相同的已知信息，策略必须给出相同动作。求解时记录当前储电量、0时计划和调整后的计划，随新观测到达再更新预测。

以题给2025年1月1日0时电量为初值，相邻日期的状态满足
\begin{equation}
  E_{\text{2025年1月1日},0}=6000,\qquad
  E_{j+1,0}=E_{j,T}.
  \label{eq:day-continuity}
\end{equation}
问题一另外要求$E_T=E_0$。问题二至四不恢复每日初值，而是接续上一日的实际末态。对于有限时长的规划，以$\Phi(E_H)$表示规划结束时剩余电量对后续费用的影响；它只参与计划选择，不作为已发生的购电支出。

\subsubsection{费用目标与近似求解方法}

设$\Pi(\mathcal I_t)$为符合信息限制和调整权限的全部策略，$K_\tau$为时段$\tau$的计划购电与调整费用。将正常支出、紧急购电及期末电量的后续影响合并，可写出期望费用目标
\begin{equation}
  \min_{\pi\in\Pi(\mathcal I_t)}
  \mathbb E\left[
    \sum_{\tau=t}^{H-1}(K_\tau+5p_\tau e_\tau)
    +\Phi(E_H)\,\middle|\,\mathcal I_t
  \right],
  \label{eq:unified-objective}
\end{equation}
求解还需满足可行域$\mathcal F$及相应初始条件。问题一不设紧急购电，日末电量必须恢复；问题二至四则通过可用信息、计划调整权限和$K_\tau$的不同规定形成各自模型。

直接求解式~\eqref{eq:unified-objective}需要同时处理所有未来情景及后续动作，计算较复杂。本文对问题一求解确定性整数规划；对问题二至四，用需求分位数形成的曲线或区间代替随机需求，建立近似线性规划，再按真实供需模拟运行。日前规划覆盖48小时，远端取$\Phi(E_{288})=0$；另用72小时及固定残值方案检验有限规划时长的影响。

评价时使用实际账单，不以近似规划的目标值代替运行费用。除总费用和期末电量外，报告日紧急费用的经验CVaR95，即最高5\%尾部的平均费用，用于反映少数高额补购日的损失。其计算方法见附录\ref{app:validation-details}。

\subsection{问题一的模型建立与求解}
\label{sec:problem1}

针对问题一，将附件一的光伏预测作为给定出力，与负荷、电价共同输入全天调度模型。先求满足日末恢复条件的最低购电费，再选取充放电总量较少、购电曲线较平稳的方案，并分析储能的节费与光伏利用效果。

\subsubsection{数据处理与购电模型}

附件一共有$144$条记录，负荷和光伏功率均非负，电价为$0.3713$--$1.3952$元/kWh。原始时间从$0{:}10$至次日$0{:}00$，结果模板则标注相应的10分钟区间。本文将时间解释为区间起点，并利用本问每日重复的条件，将末行次日$0{:}00$移至当日首段，形成$0{:}00$--$24{:}00$序列。工作簿保留原标签，指定时段按日内区间汇总；另一种时段终点解释在附录中比较。

令$\mathcal T=\{0,\ldots,T-1\}$。以计划购电量、充放电量及内部储电量为变量，建立全天费用最小的模型：
\begin{equation}
\begin{aligned}
  \min\quad & C_1=\sum_{t\in\mathcal T}p_tq_t,\\
  \text{s.t.}\quad
  &\text{式~\eqref{eq:energy-bounds}--\eqref{eq:exclusive}},\\
  &e_t=0,\quad q_t=q_t^0,\qquad t\in\mathcal T,\\
  &E_0=E_T=6000.
\end{aligned}
\label{eq:problem1}
\end{equation}
上述模型是混合整数线性规划，只保证给定预测曲线下的供电与设备可行性。将二元变量放松为连续变量，可得到最低费用的下界；若松弛解同时满足充放电互斥，则也可作为原模型的可行方案。预测出现偏差时的运行另由后续各问处理。

从经济上看，低价时购入并储存$1\ \mathrm{kWh}$，经过充放电后只能提供$\eta_c\eta_d\ \mathrm{kWh}$。若仅比较这笔电量的购入费与随后避免的购电费，低买高用能够节费的必要条件为
\begin{equation}
  p_{\rm high}\eta_c\eta_d>p_{\rm low}.
  \label{eq:arbitrage}
\end{equation}
该条件解释了价差与效率的关系，但尚未考虑容量、功率和未来光伏。比如，过早充满电池可能使午间光伏无法储存。因此，具体充放电时段仍需由全天模型联合确定。

\subsection{分层求解与最优性核验}

附件一电价严格为正，且计划未利用量不产生收益。因此，任意 $q_t>h_t$ 的可行解均可通过令 $q_t=h_t$ 降低费用，第一问求解时可以消去未利用计划变量，将普通购电与实际取电合并。保留充放电互斥变量后，模型包含 $144$ 个二元变量、$721$ 个连续变量、$288$ 个平衡及状态等式和 $288$ 个互斥线性不等式，另以变量界处理初末电量及设备边界。

使用 Python 3.9.21、SciPy 1.13.1 提供的 HiGHS 接口\cite{scipy-milp}求解，设置整数相对最优性间隙为 $10^{-10}$，每个阶段的时间上限为 $120$ 秒。首先最小化购电费，获得最优费用 $C_1^*$；随后在费用不超过 $C_1^*+10^{-7}$ 元的条件下最小化 $\sum_t(c_t+d_t)$；最后，在前两级容差内最小化日循环购电总变差
\begin{equation}
  \sum_{t=0}^{T-1}|q_t-q_{(t-1)\bmod T}|.
  \label{eq:q1-variation}
\end{equation}
后两级目标用于选择经济上近等价的方案，不能替代最低费用目标，也不保证所有时段的购电曲线平滑。最终固定已求得的充放电模式，以 $10^{-9}$ 的线性规划可行性容差重新求解连续电量，消除整数求解容差下的微小负值。

主费用整数模型报告最优解，最优性间隙为 $0$。其最优费用与线性松弛下界满足
\[
  C_1^*=C_{\rm LP}=35126.848589290\ \text{元}.
\]
线性松弛的原始目标与对偶目标一致，且其解满足充放电互斥。最终所选方案比该下界高约 $10^{-7}$ 元，处于预设的择优容差内。因此，下面的费用在分币精度上达到该模型的全局最优。

\subsection{购电与充放电结果}

主方案全天购电量为 $59482.699\ \mathrm{kWh}$，购电费为 $35126.85$ 元。指定时段购电结果见表~\ref{tab:q1-purchase}，分时段储能运行结果见表~\ref{tab:q1-storage}。完整 $144$ 时段策略保存于 \texttt{result1.xlsx}，并另存包含输入、SOC、损耗与费用的逐时段明细。

\input{../output/tables/q1_purchase}
\input{../output/tables/q1_storage}

全天母线侧充电量为 $20740.666\ \mathrm{kWh}$，放电量为 $16799.940\ \mathrm{kWh}$，初末内部储电量均为 $6000\ \mathrm{kWh}$。按式~\eqref{eq:aggregate-balance}，总负荷 $111024.808\ \mathrm{kWh}$、总光伏 $55482.836\ \mathrm{kWh}$ 与转换损耗 $3940.727\ \mathrm{kWh}$ 共同确定实际购电量；本方案光伏弃电量为 $0$。

图~\ref{fig:q1-dispatch} 展示了日内调度。凌晨与 $5{:}40$ 附近的较低价格支持储能充电，随后电池在早间负荷上升阶段放电。中午既吸收富余光伏，也在部分低价时段从外网补充电量，使储能为晚间高价时段保留资源。晚间放电后电量触及下限，$22{:}00$ 后再利用低价购电恢复日末电量。

\begin{figure}[htbp]
  \centering
  \includegraphics[width=\textwidth]{q1_dispatch.pdf}
  \caption{确定性调度的日内对应关系。a：供需与计划购电；b：放电为正、充电为负的储能功率；c：内部储电量及上下限虚线；d：电价与LP净负荷边际成本。功率与电价按10分钟区间绘制。}
  \label{fig:q1-dispatch}
\end{figure}

外网购电曲线在部分低价时段出现较高功率，其峰值为 $8458.827\ \mathrm{kW}$。这不违反题面给出的 $5000\ \mathrm{kW}$ 上限，因为后者限制储能充放电功率，而非外网总取电功率。因此，本模型的经济节省不能直接解释为外网峰值功率下降。图中的净负荷边际成本取自 LP 对偶值，用于解释给定约束状态下增加净负荷的局部代价，不作为任意幅度扰动下的固定价格。

\subsection{基线比较与独立校验}

设置无储能与分位电价规则两种基线。规则基线优先吸收剩余光伏，在已知日价格不高于 $25\%$ 分位数时充电、不低于 $75\%$ 分位数时放电，并加入期末电量可达性修正；该规则未依据优化结果调整阈值。三种策略使用相同输入与日初、日末电量条件，结果见表~\ref{tab:q1-baselines}。

\input{../output/tables/q1_baselines}

相较于无储能方案，优化方案节省 $12925.20$ 元，节费率为 $26.90\%$；相较于规则方案，购电费由 $39573.12$ 元降至 $35126.85$ 元。无储能时有 $6247.963\ \mathrm{kWh}$ 光伏无法利用，优化方案通过跨时段存储实现全部消纳。规则方案仍存在 $3656.782\ \mathrm{kWh}$ 弃光，说明仅按局部价格决定充放电不足以协调未来光伏与容量需求。

独立程序不使用优化矩阵，直接根据输出电量重新核算能量平衡、SOC 转移、功率上限、互斥和费用。主方案的能量平衡及状态转移残差均小于 $10^{-8}\ \mathrm{kWh}$，储电量保持在 $1200$--$10800\ \mathrm{kWh}$ 内，初末电量误差为 $0$。CSV 回读与官方工作簿回读检查同时通过，工作表名称、非结果单元格内容及格式保持不变。

\subsection{参数与时间解释的影响}

对容量、充放电功率及效率进行单因素重新优化，结果如图~\ref{fig:q1-sensitivity}。容量改变时同步保持 SOC 下限为额定容量的 $10\%$、上限为 $90\%$、初末电量为 $50\%$；因此，该组结果代表整套储能规格变化，不能解释为固定初始电量下的纯容量边际收益。

\begin{figure}[htbp]
  \centering
  \includegraphics[width=\textwidth]{q1_sensitivity.pdf}
  \caption{设备参数对最优费用的影响。a--c分别改变容量、功率上限与单程效率；各点均重新优化，实心点及竖虚线对应主设定。纵轴为局部费用范围，连线不表示统计置信区间。}
  \label{fig:q1-sensitivity}
\end{figure}

额定容量由 $6000$ 增至 $18000\ \mathrm{kWh}$ 时，最优费用由 $40020.40$ 降至 $32093.06$ 元；功率上限由 $2500$ 增至 $5000\ \mathrm{kW}$ 时，费用由 $36506.41$ 降至 $35126.85$ 元，再增至 $7500\ \mathrm{kW}$ 时为 $35027.13$ 元。这表明，在本日输入下，提高功率上限的边际节费趋于减弱，但该观察不能替代包含设备投资的容量选择分析。

若将题面的 $90\%$ 解释为往返效率，并取 $\eta_c=\eta_d=\sqrt{0.9}$，费用为 $33801.48$ 元，说明效率解释需要在后续各问中保持一致。光伏输入分别缩放为原值的 $0.8$ 和 $1.2$ 倍并重新优化时，费用分别为 $42268.99$ 和 $30873.96$ 元；这些属于确定性输入敏感性，不能视为既定计划应对真实预测误差的表现。

将附件时间改作时段终点，并在同样初末电量条件下重新优化，费用为 $35126.95$ 元，与主解释相差约 $0.10$ 元。虽然全天费用接近，指定时段的映射并不相同，因此仍保存两套完整策略及逐项对应关系。主结果以模板时段起点解释为准，为后续各问提供统一的时间、设备及费用核算接口。


\subsection{问题二的模型建立与求解}
\label{sec:problem2}

问题二中，实际供需不再由0时预测完全确定，而当天计划一经发布便不能修改。为处理这一变化，先分析计划过多和不足的费用差异，再利用历史误差确定购电余量，结合跨日储能安排形成计划，最后用实际供需检验运行结果。

\subsubsection{计划余量的分位数依据}
\label{sec:newsvendor}

先考虑没有储能、且计划不可调整的单个时段。设随机需求为$D\ge0$，普通电价为$p>0$，计划购电$q$全额付费，未用部分没有残值，不足部分按$5p$补购。此时期望费用为
\begin{equation}
  C(q)=pq+5p\mathbb E[(D-q)^+].
  \label{eq:single-period-cost}
\end{equation}
若需求分布连续且最优计划位于可选范围内部，对$q$求导，可得
\begin{equation}
  C'(q)=p-5p\Pr(D>q)=0
  \quad\Longrightarrow\quad F_D(q^*)=0.8.
  \label{eq:quantile}
\end{equation}
因此，在上述条件下，计划量应取需求的0.8分位数，即有80\%的概率需求不超过该值。本文据此设置多时段规划的风险参数。加入储能、后续调整或随机电价后，0.8不再必然最优，其作用是为计划余量提供一个可解释的起点；相应反例见附录\ref{app:validation-details}。

\input{../output/tables/q2_numbers.tex}
\input{../output/tables/q2_boundary_numbers.tex}

\subsubsection{供需预测与历史误差处理}

附件二包含365天、每天144个负荷与光伏记录。将时间作为时段起点，按连续日历排列：当日$0{:}00$取前一日期行的“$0{:}00{+}1$”，1月1日缺失首段以同日$0{:}10$补齐，仅用于建立初始历史。模板末列仍保留次日标签，但存放当日零点计划，作为补位例外在批注中说明；其与原标签的日期差异需结合映射表读取。另一种每日循环移位的费用差不足五百元，详见附录。\par 一月用于积累历史和另行模拟比较候选方法，实际电池约定闲置于6000 kWh。二月起连续运行334天、48096个时段，后一天继承前一天的实际电量。

在$j$日0时，根据此前历史同时预测当天和次日的负荷与光伏，记为
$\hat{\boldsymbol x}_{j,k,t}=(\hat\ell_{j,k,t},\hat v_{j,k,t})$，
其中$k=0,1$分别表示当天、次日，两类预测均只使用截至$j-1$日的数据。比较三种方法：\textbf{季节滞后法}取一周前同时段负荷与前一日光伏；\textbf{趋势法}根据最近28天的同星期负荷和最近14天的光伏变化外推，并平滑结果、缩小外推斜率及去除负值；\textbf{回归法}分别对各时段建立岭回归，直接预测两个目标日\cite{fpp-regression}：
\begin{equation}
 \hat\ell_{j+k,t}=\beta_{0,t}+\beta_{1,t}\ell_{j+k-7,t}+\beta_{2,t}\ell_{j-1,t}+\beta_{3,t}\bar\ell_{j-1},\qquad
 \hat v_{j+k,t}=\gamma_{0,t}+\gamma_{1,t}v_{j-1,t}+\gamma_{2,t}\bar v^{(7)}_{j-1,t},
 \label{eq:q2-regression}
\end{equation}
式中$\bar\ell_{j-1}$为前一日平均负荷，$\bar v^{(7)}_{j-1,t}$为最近7个完成日同时段光伏均值。回归系数用最近56个完成日重新估计，训练目标不晚于$j-1$日；标准化后岭参数取1，不足12日样本时改用趋势预测。一月统一比较预测方法及调度参数，二月起固定所选方法，回归系数仍随历史数据增加而更新。

预测当天与预测次日的误差通常不同，因此分别建立误差样本。历史上$i$日发布、目标为$i+k$日的误差为
\begin{equation}
 \boldsymbol r^{(k)}_{i,t}=\boldsymbol x_{i+k,t}-\hat{\boldsymbol x}_{i,k,t},
 \qquad i+k\le j-1.
 \label{eq:q2-matured-error}
\end{equation}
只有目标日已经结束，才能计算并使用该预测误差。例如，昨天0时对今天的预测，到今天0时仍没有完整实测值，不能加入误差样本。按这一规则随日期推进进行验证，同时检查训练数据和历史误差的可得时间\cite{fpp-tscv}。

对每类预测取最近至多56个有效历史日，将同一天的负荷、光伏误差成对加入当前预测，保留其日内顺序，并将负荷和光伏的负值截为零，得到供需情景
\[
 \boldsymbol x^{(i,k)}_{j,t}
 =[\hat{\boldsymbol x}_{j,k,t}+\boldsymbol r^{(k)}_{i,t}]^+,
 \qquad n^{(i,k)}_{j,t}=\ell^{(i,k)}_{j,t}-v^{(i,k)}_{j,t}.
\]
对各时段情景中的净负荷取经验分位数，得到用于安排购电余量的\textbf{风险需求曲线}：
\begin{equation}
 \widetilde n_{j,k,t}(\alpha)=
 \widehat F^{-1}_{j,k,t}(\alpha).
 \label{eq:q2-risk-profile}
\end{equation}
经验分位数按相邻样本线性插值计算。上述曲线逐时取分位数，0.8表示单个时段的需求分布水平，不意味着全天有80\%概率无需补购。72小时对照按相同方法增加第三日预测；年末也保持完整规划时长，跨年部分仅用已有历史外推。

\subsubsection{跨日规划与当天购电计划}

令$H=288$，将当天和次日的风险需求曲线连接，电价按已知日内轮廓重复两次。以该曲线代替随机净负荷，沿用电量平衡和设备约束，建立近似规划：
\begin{equation}
\begin{aligned}
 \min\quad &\sum_{\tau=0}^{H-1}p_\tau\widetilde q_\tau+\Phi_H(E_H),\\
 \text{s.t.}\quad&
 \widetilde q_\tau+d_\tau=\widetilde n_\tau+c_\tau+s_\tau,\quad
 \widetilde q_\tau,s_\tau\ge0,\\
 &E_{\tau+1}=E_\tau+\eta_cc_\tau-d_\tau/\eta_d,\\
 &E^{\min}\le E_\tau\le E^{\max},\quad
 0\le c_\tau\le\bar c,\quad0\le d_\tau\le\bar d.
\end{aligned}
\label{eq:q2-proxy}
\end{equation}
初始储电量取当天实际值，$s_\tau$表示规划中未利用的余电，主方案不计两天后电量的后续价值，即$\Phi_H=0$。求解后仅把前144个$\widetilde q_\tau$确定为当天计划；次日部分只帮助判断应保留多少储能，届时仍重新制定计划\cite{mpc-horizon}。

由于当天24时处于两天规划的中间，其储电量会影响次日购电费。若当天末值为$e$，次日最低规划费用可表示为
\begin{equation}
 J^{\rm next}_{j}(e)=
 \min_{\text{次日可行轨迹},\,E_{144}=e}
 \sum_{\tau=144}^{287}p_\tau\widetilde q_\tau+\Phi_{288}(E_{288}).
 \label{eq:q2-continuation}
\end{equation}
式~\eqref{eq:q2-continuation}说明，当天末值的作用由次日需求、电价和设备条件共同决定。两天之后的价值仍被截断，因此还需与72小时规划及计入固定残值的方案比较，判断截断是否影响当天决策。

用HiGHS求解线性规划，并检查约束残差与原始、对偶目标差\cite{scipy-linprog}。若解中出现同时充放电，令$\varepsilon=\min\{c,d/(\eta_c\eta_d)\}$，将$(c,d)$改为$(c-\varepsilon,d-\eta_c\eta_d\varepsilon)$；这样内部储电量变化不变，在正电价且余电允许弃用时，所需购电不会增加。规划完成后，还要根据实际供需逐段执行，不能把预测曲线下的可行性当作实际无缺口的保证。

\subsubsection{实际充放电规则与参数选择}

0时确定计划后，电池负责吸收计划与实际供需的偏差。设当前余缺电量为$b_t=q_t+v_t-\ell_t$，采用“有余电则充电、有缺口则放电”的贪心规则：
\begin{equation}
\begin{aligned}
 b_t\ge0:&\quad c_t=\min\{b_t,\bar c,(E^{\max}-E_t)/\eta_c\},\quad d_t=0;\\
 b_t<0:&\quad d_t=\min\{-b_t,\bar d,(E_t-E^{\min})\eta_d\},\quad c_t=0.
\end{aligned}
\label{eq:q2-greedy}
\end{equation}
式~\eqref{eq:q2-greedy}中，充放电量同时受余缺电量、功率和内部电量限制。运行时优先使用光伏，再取用已付费的计划电量，电池放电仍不足的部分按5倍电价补购。该规则容易执行，但提前放电可能耗尽晚间高价时段的储量，因此不保证费用最优。\par 回放以当前10分钟的总电量计算动作，近似设备在时段内的快速反馈。若要求时段开始前一次性固定充放电量，当前完整实测值便不可用，需改为预测驱动并另行检查控制延迟。本文运行结果以前述快速反馈近似为条件。

另设两种充放电规则作对照：一是跟踪0时规划给出的参考储电量，二是每小时利用已观测误差更新预测，再优化剩余储能动作。三者均不修改当天购电计划，并连续传递实际电量。

一月15--31日对三种预测方法、四个分位数$\{0.50,0.65,0.80,0.90\}$及三种充放电规则分别模拟，共36组。按计入共同剩余电量价值后的费用排序，差额不超过0.01元视为并列，再按预设规则选择，最终采用\QBForecast、$\alpha=\QBQuantile$及上述贪心规则，验证费用为\QBValidation 元，并列\QBTies 项。二月起保持选型不变，历史样本逐日更新。48小时结构曾依据全年诊断修订，因此后续评价属于按信息顺序进行的历史回放。

\subsubsection{购电费用与紧急购电风险}

连续运行所得计划费为\QTwoPlanCost 万元，紧急费为\QTwoEmergencyCost 万元，合计\QTwoTotalCost 万元；紧急购电\QTwoEmergencyEnergy\ kWh，期末储电量\QTwoTerminal\ kWh。不同策略的期末电量并不相同，为避免将多耗电池电量误认为节费，另计算统一修正费用
\begin{equation}
 \nu_{\rm eval}=\min_t p_t/\eta_c=\QBCommonNu,\qquad
 C^{\rm adj}=C^{\rm cash}-\nu_{\rm eval}(E_{\rm end}-E_{\rm start}).
\end{equation}
其中$\nu_{\rm eval}$只用于结果比较，不能与规划中的期末价值混用。主方案的修正费用为\QBAdjusted 元。实际运行从2月1日6000 kWh开始，当天末值为\QBFirstEnd\ kWh，此后各日接续运行。

\input{../output/tables/q2_comparison.tex}

无储能对照采用相同预测和0.8分位数，点预测对照保留相同规划时长与充放电规则，但不加历史误差形成的余量。主方案相对两者分别节费\QTwoSaving\%和\QBPointSaving\%。图~\ref{fig:q2-cost}a显示需求曲线、计划与实际补购的对应关系；b表明，预留较多计划电量虽增加正常购电费，却减少了更多高价补购支出。

\begin{figure}[htbp]
 \centering\includegraphics[width=\textwidth]{../output/figures/q2_risk_planning.pdf}
 \caption{问题二的计划余量与费用比较。a：指定日期中紧急购电最多的9月23日，灰带为历史误差情景的净负荷分布范围，实线为0.8分位数需求曲线，蓝色为0时计划，斜纹为补购；b：334天计划费与紧急费；c：主方案相对各对照的日均节省及七日块重采样95\%区间；d：总费用与日紧急费用CVaR95。}
 \label{fig:q2-cost}
\end{figure}

当天预测的负荷、光伏平均绝对误差（MAE）分别为\QTwoLoadMAE\ kW和\QTwoPVMAE\ kW；实际净负荷不超过风险需求曲线的时段占\QTwoCoverage\%。误差仅在已有实测的目标时段计算，不对年末缺乏实测的辅助预测评价精度。

主方案有\QBEmergencyDays 天、\QBEmergencySlots 个时段需要紧急购电，补购后负荷均被满足。日紧急费用CVaR95为\QBCvar 元，无储能对照为\QBBaseCvar 元。总费用与少数高损失日的表现应分别评价，不能仅因平均节费就认定风险全面降低；图~\ref{fig:q2-cost}c、d给出两方面结果，区间的解释见\hyperref[sec:model-validation]{模型检验与评价}一节。

\subsubsection{规划时长与期末电量的影响}

为检验跨日规划的作用，固定预测、分位数、充放电规则及初态，比较五种设置：24小时原终端规则、24小时替代成本方案A、48小时方案B、72小时方案及48小时计入远端残值的方案。A取$\nu=p_{0:00}/\eta_c$，远端残值取$\nu=\min p/\eta_c$，结果均用$\nu_{\rm eval}$统一修正。表中另外保留旧配置结果作为历史参照。

\input{../output/tables/q2_boundary.tex}

旧配置采用两侧各0.9的效率、每日循环时标和原预测方法集合，与当前配置不同，不能作为只改变规划时长的对照。以下只比较相同参数的结果：B与24小时方案的计划费差为\QBMatchedPlanDifference 元（B减24小时），紧急费差为\QBMatchedEmergencyDifference 元（24小时减B），现金费用相对变化\QBMatchedPercent\%；期末电量修正后的结果见表~\ref{tab:q2-boundary}。

与B相比，72小时方案和加入远端残值方案的费用变化分别为\QBSeventyTwoDifference 元、\QBSalvageDifference 元，计划和执行记录也一致，表明这两项变化未改变本次当天决策。B相对24小时方案的日均修正节费95\%区间为$[\QBCIMatchedLow,\QBCIMatchedHigh]$元，包含零，因而尚不能确认48小时方案具有稳定费用优势。

附录图~\ref{fig:q2-boundary}进一步给出晚间储能变化。日末平均电量由\QBOldEndMean 升至\QBEndMean\ kWh，22--24时日均充电由\QBOldLateCharge 变为\QBLateCharge\ kWh，该时段累计紧急购电由\QBOldLateEmergency 变为\QBLateEmergency\ kWh。延长规划主要改变了晚间保留和补充电量的方式，新增计划费与补购费变化相互抵消；晚间补购还受预测误差及设备限制影响，不能全部归于期末处理。

\subsubsection{指定日期结果与参数敏感性}

指定日期的0时计划和实际总费用见表~\ref{tab:q2-purchase}，充放电及初末电量见表~\ref{tab:q2-storage}，紧急购电区间见表~\ref{tab:q2-emergency}。这些日期均取自连续运行记录，未单独重置电池。完整逐时曲线见附录图~\ref{fig:q2-days}。

\begingroup
\renewcommand{\arraystretch}{0.9}
\input{../output/tables/q2_purchase_layout.tex}
\input{../output/tables/q2_storage_layout.tex}
\endgroup
\input{../output/tables/q2_emergency.tex}

表~\ref{tab:q2-emergency}中的“--”表示当天没有紧急购电，合计0.000不表示没有向外网购电。3月20日、6月21日和12月21日的计划缺口由电池补足；9月23日20:20先受放电功率限制，20:30--20:40又降至电量下限，需要继续补购至22时。四日由题目指定，全年仍有\QBEmergencyDays 个紧急购电日，不能据此判断其他日期也无风险。

另对指定日完成\QBSensitivityCases 组敏感性计算，分别改变分位数、设备参数、补购倍率和误差幅度。每组从当天实际初态重新运行，容量变化时保持初始储电比例，误差变化仅影响预测而不改动真实供需。这些结果反映四日局部响应；规划时长和残值则采用全评价期对照。逐时记录及工作簿检查见\hyperref[sec:model-validation]{模型检验与评价}一节和附录\ref{app:validation-details}。


\subsection{问题三的模型建立与求解}
\label{sec:problem3}

与问题二相比，本问允许在新预报到达后修改购电计划。此时既要利用更新后的供需信息，也要从电池的实际电量出发计算剩余供能能力。因此，先确定增加和取消计划的收费方式，再建立各调整时点的剩余时段优化模型。

\subsubsection{计划更新与调整费用}

以$\mathcal K=\{0,36,72,108\}$表示每日0、6、12、18时对应的时段序号。设$q_t^{(k)}$为在时点$k$确定的时段$t$购电计划，$q_t^{(0)}=q_t^0$。当$k>0$时，仅修改$t\ge k$的未执行部分，已经供电的时段不再改变。

每次使用截至当前最新发布的光伏预报，将整点值插值到10分钟。预报覆盖次日的部分用于判断后续储能需求，不作为次日正式购电计划提前执行。

设时段$t$最后生效的购电量为$q_t$，与0时计划相比，其增加量和取消量分别为
\begin{equation}
  \delta_t^+=(q_t-q_t^0)^+,\qquad
  \delta_t^-=(q_t^0-q_t)^+.
  \label{eq:adjustment}
\end{equation}
对于取消部分，题面没有唯一说明原购电费是否退还。本文采用退回一半原费用的解释：保留电量按原价付费，增加电量按1.5倍付费，取消电量退回0.5倍原价。因此，计划与调整费用为
\begin{equation}
\begin{aligned}
  K_t&=p_t\min(q_t^0,q_t)+1.5p_t\delta_t^++0.5p_t\delta_t^-\\
     &=p_tq_t^0+1.5p_t\delta_t^+-0.5p_t\delta_t^-
      =p_tq_t+0.5p_t\lvert q_t-q_t^0\rvert.
\end{aligned}
  \label{eq:contract-cost}
\end{equation}
上式也可理解为：最终计划按原价结算，再对偏离0时计划的电量加收50\%费用。加上实际紧急购电后，总现金支出为
\begin{equation}
  C_3^{\rm cash}=\sum_{t\in\mathcal T}
  \left(p_tq_t^0+1.5p_t\delta_t^+-0.5p_t\delta_t^-+5p_te_t\right).
  \label{eq:problem3-cash}
\end{equation}
本文按最终计划与0时计划的净差结算一次。若另一种规则要求每次相邻版本的变化都收费，则先增后减等往返调整也会产生费用，不能用最终净差抵销；此时模型还需继承上一版有效计划。后文结果均以本节的净额结算约定为前提。

另设不退款的解释作为对照：0时计划全额付费，取消后不退回原费用，并加收50\%违约费。其计划与调整费用为
\begin{equation}
  K_t^{\rm pen}=p_tq_t^0+1.5p_t\delta_t^++0.5p_t\delta_t^-.
  \label{eq:penalty-cost}
\end{equation}
电价为正、闲置计划允许不取用时，不退款规则下主动取消电量只会增加费用；退款规则则可收回部分支出。两种解释分别运行的结果见第\ref{sec:q3-settlement}节。

明确费用后，在6、12、18时根据当前预报和实际储电量重算余下计划。下面先说明外部预报及历史误差的处理，再给出调整模型与运行结果。

\input{../output/tables/q3_numbers.tex}

\subsection{发布预报的接入与融合}

附件三在每日$0{:}00$、$6{:}00$、$12{:}00$、$18{:}00$给出其后$1$至$24$小时的整点光伏功率。将“预报$h$小时”分别对齐到发布后$h$小时起始的10分钟时段、其前一时段、其前一小时均值和其后一小时均值，并与附件二实测比较（全年\QThreeAlignDays 个有完整实测的日期，归档文件\texttt{alignment\_check.csv}），$0{:}00$版本的平均绝对误差分别为\QThreeAlignStart、\QThreeAlignEnd、\QThreeAlignPrevMean、\QThreeAlignNextMean\ kW，其余三个版本的排序相同。因此本文将整点预报锚定于发布后$h$小时起始的时段。发布前刚结束时段的实测值是发布时刻已知的量，作为首个锚点；锚点之间线性插值到10分钟并截断为非负。发布之后的实测不进入任何版本的转换，测试通过改写发布后的观测验证这一点。发布时刻$b$的版本覆盖当日$b$至$143$与次日$0$至$b-1$共144个时段，即提前期$1$至$144$个时段；下文按“发布时刻—目标时段—提前期”明确每个版本的使用范围。

外部预报并不总优于本地预测。表~\ref{tab:q3-forecast}表明，发布后24小时窗口内，$0{:}00$版本的误差高于第二问的本地趋势预测，$6{:}00$版本低于本地预测，$12{:}00$和$18{:}00$版本因跨越次日而误差更大。两者等权融合在$0{:}00$、$6{:}00$、$12{:}00$均低于任一单独来源；该权重取$0.5$，是在查看这组误差诊断之后、任何费用回放之前固定的，此后不按费用结果调整（本文所称“先验”均指相对费用回放而言）。对当日剩余时段而言，最新版本明显优于$0{:}00$版本：$6{:}00$时由\QThreeStaleMAESix\ kW降到\QThreeFreshMAESix\ kW，$12{:}00$时由\QThreeStaleMAETwelve\ kW降到\QThreeFreshMAETwelve\ kW；$18{:}00$后当日已无光伏，该版本只能通过次日辅助段起作用。对各版本覆盖的次日部分（$6{:}00$、$12{:}00$、$18{:}00$版本分别对应提前约18至24、12至24、6至24小时），等权融合反而不如本地两日提前期预测：$12{:}00$版本为\QThreeNextAwareTwelve\ kW对\QThreeNextLocalTwelve\ kW，$18{:}00$版本为\QThreeNextAwareEighteen\ kW对\QThreeNextLocalEighteen\ kW。本文按先验保持单一权重，把按提前期设定权重（只用已成熟误差选取）留作后续工作。

\input{../output/tables/q3_forecast.tex}

主模型的光伏点预报取融合值，负荷沿用第二问的趋势预测；在调整时刻，负荷点预报按最近6个时段的实测误差均值作时间常数为12个时段的指数衰减修正。风险情景与第二问同构：对每个发布版本，用同一融合方式重建历史各日的预报，只取目标日已结束的成对残差（最近56日），叠加到当前点预报后取净负荷分位数；历史负荷预测同样重建各历史日期在同一调整时刻本可执行的偏差修正，使残差库与实际使用的预测器同构。次日辅助段在调整时把当前版本的次日覆盖部分与本地两日提前期预测融合，覆盖之外的次日时段沿用本地预测，不读取次日才发布的版本；其残差按同一构造、只用成熟目标重建。以上两处为第一轮审计后补全的信息链；此前版本的次日辅助段沿用$0{:}00$本地预测且负荷残差未作同构重建，本文保留为“冻结旧版”对照，并在第\ref{sec:q3-chain}节同参比较。

\subsection{有偿调整的决策模型}

结算按式~\eqref{eq:contract-cost}：每个时段以最后有效合同$q_t$与$0{:}00$计划$q_t^0$一次结算，上调部分付$1.5p_t$，取消部分退回$0.5p_t$。由此得到三个边际获取价格：日前计划$p_t$、日内上调$1.5p_t$、放弃取消即保留计划$0.5p_t$。按第\ref{sec:newsvendor}节的推导，在单时段、无储能情形下对应的最优服务分位数分别为
\begin{equation}
  \alpha^{0}=1-\frac{1}{5}=0.8,\qquad
  \alpha^{+}=1-\frac{1.5}{5}=0.7,\qquad
  \alpha^{-}=1-\frac{0.5}{5}=0.9.
  \label{eq:q3-band}
\end{equation}
即计划低于更新分布的$0.7$分位数时才值得上调，高于$0.9$分位数时才值得取消，两者之间保持计划不动。$[\alpha^{+},\alpha^{-}]$称为调整带，其两个端点由题面费用比例导出；$0{:}00$沿用$0.8$分位数则是继承第二问的近似，第\ref{sec:newsvendor}节已说明它对含后续调整机会的机制并非最优。

设调整时刻为$k$，当日剩余时段为$t=k,\dots,143$，次日辅助段为$t=144,\dots,287$。记$\widetilde n_t^{-}=\widetilde n_t(\alpha^{+})$、$\widetilde n_t^{+}=\widetilde n_t(\alpha^{-})$，调整LP为
\begin{equation}
\begin{aligned}
 \min\quad&\sum_{t=k}^{143}p_t\bigl(1.5u_t-0.5v_t-w_t\bigr)+\sum_{t=144}^{287}p_t\widetilde q_t,\\
 \text{s.t.}\quad&q_t^0+u_t-v_t+d_t=\widetilde n_t^{-}+w_t+c_t+s_t,\quad t<144,\\
 &0\le v_t\le q_t^0,\quad 0\le w_t\le\widetilde n_t^{+}-\widetilde n_t^{-},\quad t<144,\\
 &\widetilde q_t+d_t=\widetilde n_t(\alpha^0)+c_t+s_t,\quad \widetilde q_t\ge0,\quad t\ge144,\\
 &E_{t+1}=E_t+\eta_cc_t-d_t/\eta_d,\quad E_k=E_k^{\rm actual},\\
 &E^{\min}\le E_t\le E^{\max},\quad 0\le c_t\le\bar c,\quad 0\le d_t\le\bar d.
\end{aligned}
\label{eq:q3-adjust}
\end{equation}
其中$u_t,v_t$为相对$0{:}00$计划的上调与取消量，$w_t$为调整带内被覆盖的需求，其单位价值取计划价$p_t$：覆盖带内需求只在代价低于计划价时发生，因此不会以$1.5p_t$增购去覆盖它，也不会为$0.5p_t$的退款放弃它，在单时段、无储能情形下恰好实现式~\eqref{eq:q3-band}的规则。带内单位价值取常数$p_t$是一个额外的分段线性代理假设：任何介于$0.5p_t$与$1.5p_t$之间的常数都给出同样的单时段不动区间，但有储能时会改变参考轨迹（附录测试给出两时段算例），真实补购的边际价值则随覆盖水平与条件尾部概率变化。因此该LP是保持单时段经济调整带的代理规划，不是原随机控制问题的等价重写。储能以实际$E_k$为初态参与权衡。次日变量只用于给出日末电量价值，不签约、不结算；远端终端值与第二问一致取零，视窗末端固定为次日$24{:}00$，故$6{:}00$、$12{:}00$、$18{:}00$的剩余视窗分别为42、36、30小时。每次求解后只提交$t<144$的新合同，$t<k$的合同和费用不变。线性规划用HiGHS求解并核验对偶证书，同时充放电按第二问方法消除。

\subsection{因果执行与一月检验}

每日$0{:}00$以实际电量为初态、按第二问的48小时风险规划提交当天合同，其中当日光伏点预报改为融合预报。在$6{:}00$、$12{:}00$、$18{:}00$，决策只使用该时刻的实际SOC、此前时段的实测负荷与光伏、当时已发布的最新版本（含其次日覆盖部分）以及$0{:}00$计划。执行采用第二问选定的贪心缓冲控制器，普通合同按最后有效版本取电，紧急电只补缺口。主策略先验固定为“三个时点全部允许调整、退款口径、调整带$[0.7,0.9]$、融合权重$0.5$”，不据全年结果改选；信息链修复按事先约定的规则采纳为正式版本：正式发布的版本必须是其声明的信息集与实现一致的版本，费用差异只作报告。

一月15日至31日对调整规则作影子检验，各变体自6000 kWh出发、按统一残值修正：导出带$[0.7,0.9]$为\QThreeCheckBand 元，单一$0.8$分位数为\QThreeCheckSymmetric 元，宽带$[0.6,0.95]$为\QThreeCheckWide 元，不调整为\QThreeCheckPlanOnly 元，冻结旧版信息链为\QThreeCheckLegacy 元。导出带略优于替代规则，与先验一致；该检验只作报告，不用于选型。

\subsection{全年结果与指定日期}

二月至十二月，主策略的计划费为\QThreePlanCost 万元，调整结算为\QThreeAdjustCost 万元（上调\QThreeUpCost 万元，取消退回\QThreeDownCost 万元），紧急费为\QThreeEmergencyCost 万元，现金总费用为\QThreeTotalCost 万元。全年上调\QThreeUpEnergy\ kWh、取消\QThreeDownEnergy\ kWh，在\QThreeAdjustedDays 天的\QThreeAdjustedSlots 个时段改变了合同；紧急购电\QThreeEmergencyEnergy\ kWh，发生于\QThreeEmergencyDays 天，日紧急费用经验CVaR95为\QThreeCvar 元。期末电量\QThreeTerminal\ kWh，统一残值修正费用\QThreeAdjusted 元。

相对同一融合预报、仅$0{:}00$计划的策略，现金节省\QThreeSavingPlanOnly 元，节费率\QThreeSavingPlanOnlyPercent\%，紧急购电由\QThreePlanOnlyEmergency\ kWh降至\QThreeEmergencyEnergy\ kWh，CVaR95由\QThreePlanOnlyCvar 元降至\QThreeCvar 元；配对七日块重采样的日均节省95\%区间为$[\QThreeCIPlanOnlyLow,\QThreeCIPlanOnlyHigh]$元，不含零。相对第二问已发布策略（不用附件三、不调整），节省\QThreeSavingQTwo 元，节费率\QThreeSavingQTwoPercent\%。取消量远大于上调量，说明$0{:}00$计划在$0.8$分位数上保留的余量在信息更新后大部分被有偿收回，而上调集中在少数出现缺口的时段。

指定日期的结果见表~\ref{tab:q3-purchase}至表~\ref{tab:q3-emergency}：表~\ref{tab:q3-purchase}每格同时给出$0{:}00$计划量与最终合同量，全天购电费为含调整结算与紧急费的实际账单；表~\ref{tab:q3-storage}为最终执行的充放电量与日初日末电量；表~\ref{tab:q3-emergency}为最终紧急购电区间。附录图~\ref{fig:q3-days}显示，调整主要发生在两类时段：中午光伏高于计划预期时取消合同，傍晚$18{:}00$后依据实际电量与负荷水平上调合同，使晚高峰的紧急购电基本消失。

\begingroup
\renewcommand{\arraystretch}{0.9}
\input{../output/tables/q3_purchase.tex}
\input{../output/tables/q3_storage.tex}
\endgroup
\input{../output/tables/q3_emergency.tex}

\subsection{信息链修复的同参对照}
\label{sec:q3-chain}

表~\ref{tab:q3-information-chain}在相同输入、初态、控制器与结算下比较冻结旧版、仅接入跨日新预报、仅同构残差校准与两项同时修复（主策略）。冻结旧版费用为\QThreeLegacyCost 元，主策略比它少\QThreeLegacyDiff 元，日均差的95\%区间$[\QThreeCILegacyLow,\QThreeCILegacyHigh]$元含零；差额几乎全部来自同构残差（仅同构变体与主策略相差\QThreeMatchedBiasDiff 元），仅接入跨日新预报的变体与旧版仅相差\QThreeCrossDayVsLegacy 元。机制上，同构重建后$18{:}00$调整带的实际覆盖率由旧版的\QThreeLegacyCovEighteenTodayLow\%/\QThreeLegacyCovEighteenTodayHigh\%回到\QThreeCovEighteenTodayLow\%/\QThreeCovEighteenTodayHigh\%（名义$70\%/90\%$），说明旧残差库叠加完整日前误差后在傍晚偏保守；带宽收窄使上调减少（调整结算减少\QThreeLegacyAdjustDiff 元），紧急购电由\QThreeLegacyEmergency\ kWh增至\QThreeEmergencyEnergy\ kWh，CVaR95由\QThreeLegacyCvar 元增至\QThreeCvar 元。跨日新预报影响很小的原因有二：在夜间便宜、晚高峰昂贵的价格结构下，今晚的充电合同由次日晚高峰决定，次日午间光伏信息只改变白天的储能参考轨迹；且各版本次日覆盖部分的光伏误差高于本地两日预测，$0.8$分位数轮廓在同构残差下吸收了这一差异（次日覆盖段的实际覆盖率为\QThreeCovEighteenCoveredLow\%）。因此这两项修复是信息链的逻辑修正，而非策略性能改进；它们被采纳为正式版本的依据是信息集声明与实现一致，不是费用更低，尾部风险略升的负面结果一并报告。

\input{../output/tables/q3_information_chain.tex}

\subsection{信息价值分解与预报时点的必要性}

表~\ref{tab:q3-comparison}在相同实际轨迹、初态、执行器和结算口径下比较各策略。按式~\eqref{eq:voi}逐层分解（图~\ref{fig:q3-value}a）：从第二问策略到“融合预报、仅计划”，$0{:}00$外部预报经融合进入计划带来\QThreeValueExternalPlan 元；从“仅计划”到“三时点调整但光伏预报冻结为$0{:}00$版本”节省\QThreeValueAdjustNoInfo 元，这一步同时包含合同调整权限与实际SOC、负荷偏差的状态反馈，现有对照不能再把二者分开；再改用最新发布版本（含当日剩余与次日覆盖部分），只再节省\QThreeValueFreshInfo 元，日均差的95\%区间$[\QThreeCIStaleLow,\QThreeCIStaleHigh]$元跨过零。因此，调整机制的主要价值来自合同调整权限与状态反馈，而不是日内光伏预报精度的提高；后者的作用与图~\ref{fig:q3-value}d所示的误差下降并不成比例，原因是白天光伏偏差多数可由储能吸收，而费用集中的晚高峰缺口与光伏预报无关。

\input{../output/tables/q3_comparison.tex}

各时点的边际价值互不相同（图~\ref{fig:q3-value}b）。在三时点基础上去掉$6{:}00$、$12{:}00$、$18{:}00$，费用分别增加\QThreeMarginalSix、\QThreeMarginalTwelve、\QThreeMarginalEighteen 元；单独使用一个时点时，节省分别为\QThreeSingleSix、\QThreeSingleTwelve、\QThreeSingleEighteen 元。$18{:}00$最为重要：此时晚高峰前的实际电量、当日负荷水平已经确定，恰可用$1.5$倍电价的上调替代$5$倍电价的紧急补购；这是交易权限与状态反馈的价值，不能全部归于$18{:}00$光伏预报本身。$12{:}00$次之，其作用主要是收回中午被高估的计划电量（图~\ref{fig:q3-value}c）。$6{:}00$的边际价值最小，去掉它的日均费用差区间$[\QThreeCIDropSixLow,\QThreeCIDropSixHigh]$元跨过零。就“是否需要引入其他时刻的预报”而言：三个时点的边际收益均为正，全部引入是主策略；若只能保留一个，应保留$18{:}00$；若保留两个，应为$12{:}00$与$18{:}00$。附件之外并不存在其他时刻的真实预报，本文不虚构新增版本；上述结果表明，若要增设时点，优先级应按晚高峰前的状态反馈价值而非光伏预报精度排序，而评价新增预报服务本身还需固定交易权限另作对照。

\begin{figure}[htbp]
 \centering\includegraphics[width=\textwidth]{../output/figures/q3_information_value.pdf}
 \caption{信息价值与预报误差。a：四个策略的334天现金费用阶梯，标注相邻差值，竖线为各策略与主策略日均费用差的95\%配对区间（换算为334天）；b：单一时点相对仅计划策略的节省，以及从三时点中去掉该时点的损失，误差线为95\%配对区间；c：各发布版本按目标小时的日均上调与取消电量，取消为负并加斜纹；d：各发布版本24小时窗口的光伏MAE。}
 \label{fig:q3-value}
\end{figure}

\begin{sloppypar}
外部预报的作用需要与融合一起理解。仅用附件三预报的三时点策略费用为 \QThreeExternalOnlyCost 元，比主策略高 \QThreeExternalOnlyDiff 元。完全不用附件三但允许三时点调整的费用为 \QThreeOwnAdjustCost 元，而主策略为 \QThreeCash 元。因此附件三的价值主要通过与本地预测融合实现，单独依赖外部预报反而不如融合，而不用外部预报时调整机制本身仍保留大部分收益。第二问策略在第三问账本下的费用与第二问发布结果完全一致，验证了两问账本的连续性。
\end{sloppypar}

\subsection{结算口径与调整规则的稳健性}
\label{sec:q3-settlement}

在不退款口径式~\eqref{eq:penalty-cost}下，同一策略的取消量为\QThreePenaltyDown\ kWh，费用为\QThreePenaltyCost 元，比退款口径高\QThreePenaltyDiff 元：取消机制被架空后，计划余量只能作为沉没成本，节省完全来自上调。这一差异是结算规则的直接后果，本文以字面读法即退款口径为主口径，并如实报告两者。

将调整带改为单一$0.8$分位数，费用为\QThreeSymmetricCost 元，与主策略之差为\QThreeSymmetricDiff 元（负值表示对称规则更便宜），日均差区间$[\QThreeCISymmetricLow,\QThreeCISymmetricHigh]$元，一月影子检验则略偏向导出带。该对照在开发过程中已被反复查看，不能据其追溯改选主策略；它与第\ref{sec:q3-chain}节的发现一致地指出，代理调整带的带内价值函数与真实补购损失之间仍有差距，这是下一阶段修正的对象。保留导出带是因为其端点由费用比例推出、无需校准，而非因为回放费用最低。融合权重、残差窗口与负荷偏差修正参数均为先验设定，未作全年调参。

\subsection{独立核验}

独立程序对\QThreeStrategyCount 条全年轨迹重算能量平衡、SOC转移与连续性、设备边界、充放电互斥、紧急电用途、合同取电恒等式、调整恒等式$q_t=q_t^0+u_t-v_t$及$u_tv_t=0$、合同版本不早于其生效时段、首个调整时刻前无调整，以及计划费、调整结算、紧急费与总费用的逐项恒等式，最大残差不超过$2\times10^{-5}$ kWh。全年\QThreeCash 元的现金费用由\QThreePlanCost、\QThreeAdjustCost、\QThreeEmergencyCost 万元三项精确加总。这些轨迹共求解\QThreeLPSolves 个线性规划（含一月影子检验全程为\QThreeLPSolvesTotal 个），最大原始--对偶目标差$\QThreeDualGap$元，被状态、残差或对偶证书拒绝的求解为\QThreeSolverFailures 次（任一拒绝都会中止回放，故完整归档的该计数必为零）。信息链四组变体的$0{:}00$计划288个时段与每份调整记录（主策略\QThreeAdjustmentRecords 次）的完整辅助段均已保存：辅助段满足状态转移、电量与功率边界，标记为未承诺，不进入任何账单；其点预报、风险轮廓与偏差修正由原始附件在归档配置下因果重建，最大差在$10^{-5}$ kWh以内，并对抽样记录在改写发布后观测、后续日期与未发布版本后重建第二次，结果逐项相同。测试用例改写尚未发布的版本与发布后的实测，确认此前的计划、合同与执行不变；两个信息链开关均关闭时逐时复现2026年9月11日的冻结归档（最大差$10^{-12}$量级）；权重取零且不调整时逐时复现第二问已发布策略。官方工作簿\texttt{result3.xlsx}的两张购电表分别为$0{:}00$计划及其费用、最终有效合同及全天实际账单，充放电量与紧急购电量按最终执行填写，回读与轨迹一致。


\subsection{问题四的模型建立与求解}
\label{sec:problem4}

问题四把前两问的固定分时价格改为随日期变化的实时价格，未来实价在决策时未知。因此，需要增加电价预测，分别重新计算全天计划不变和日内可调整两种策略。本节先分析价格变化，再建立预测方法，最后比较结算、预测和调整各自引起的费用差异。

\input{../output/tables/q4_numbers.tex}

\subsection{价格结构与因果预测}

附件四价格范围为0.008--1.794元/kWh，全年均值与附件一相同，平均日内轮廓相关系数为1.0；主要变化体现在日级水平和日内偏差（图~\ref{fig:q4-prices}）。逐时价格与净负荷的相关系数为\QFourPriceNetCorr，逐时段回归$R^2$为\QFourPriceRSquared。净负荷预测误差最大的20\%时段，电价为均值的\QFourDeficitPriceRatio 倍，提示补购风险与高价可能叠加。

据此，以最近两个同星期日的价格轮廓与最近56日逐时段净负荷回归等权预测价格，回归输入采用本地净负荷预测，两个提前期共用已完成日期。组合由一月8--31日误差检验确定，二月起冻结。6、12、18时再以当日已观测的平均价格偏差修正剩余预测，形成“日前轮廓＋日内偏差”的更新方式。

\input{../output/tables/q4_price.tex}

二月至十二月，融合预测MAE为\QFourMAECombined 元/kWh，低于固定轮廓的\QFourMAEFixed 元/kWh；三次日内修正也降低了剩余时段MAE。预测是否节费，仍需在相同结算规则下通过运行对照判断。

\begin{figure}[htbp]
 \centering\includegraphics[width=\textwidth]{../output/figures/q4_prices.pdf}
 \caption{实时电价的结构与可预测性。a：附件四实时电价按星期与10分钟时段的二月至十二月均值；b：各日平均电价与平均净负荷，周五、周六由价格聚类识别，虚线为最小二乘拟合；c：二月至十二月各因果日前预测方法的当日平均绝对误差，条形从零起。}
 \label{fig:q4-prices}
\end{figure}

\FloatBarrier
\subsection{两类计划调整机制与结算模型}
\label{sec:q4-price-convention}

4-2沿用本地预测并保持0时计划购电量不变，4-3沿用融合光伏预报、三时点调整及其信息链。规划时以预测价格替换相应代理模型中的$p_t$，实际现金费用则按相应时段实价计算：
\begin{equation}
  C_{4\text{-}2}^{\rm cash}=\sum_t p_t(q_t^0+5e_t),
  \label{eq:problem42}
\end{equation}
\begin{equation}
  C_{4\text{-}3}^{\rm cash}
  =\sum_t p_t(q_t^0+1.5\delta_t^+-0.5\delta_t^-+5e_t).
  \label{eq:problem43}
\end{equation}
实际价格只进入结算和已到达的观测，不提前供规划使用。两族保持设备、执行器和残差构造一致，并分别设置固定轮廓及当日价格预知对照；4-3另设无日内价格修正和仅0时计划。价格预知只替换输入价格，保留代理目标与执行器，因而不构成最优费用界限。

当价格$P$与需求$D$相关时，单时段最优条件为$\mathbb E[P]=5\mathbb E[P\mathbf1_{\{D>q\}}]$，普通分位数不再自动适用。本扩展暂继承0.8分位数与$[0.7,0.9]$调整带以保持机制可比，其结论属于因果点价格预测下的代理调度。

\subsection{运行结果与调整收益}

\textbf{计划固定（4-2）。}现金费用为\QFourTwoTotalCost 万元，其中计划费\QFourTwoPlanCost 万元、紧急费\QFourTwoEmergencyCost 万元；紧急购电\QFourTwoEmergencyEnergy\ kWh，期末电量\QFourTwoTerminal\ kWh，日紧急费用CVaR95为\QFourTwoCvar 元。相对固定电价下的问题二，费用增加\QFourTwoVsQTwoPercent\%，说明在相同平均电价下，价格与购电时序的匹配仍会改变支出。

\textbf{有偿调整（4-3）。}现金费用为\QFourThreeTotalCost 万元，由计划费\QFourThreePlanCost 万元、调整结算\QFourThreeAdjustCost 万元与紧急费\QFourThreeEmergencyCost 万元组成。紧急购电\QFourThreeEmergencyEnergy\ kWh，CVaR95为\QFourThreeCvar 元。相对同族仅0时计划，节省\QFourThreeSavingPlanOnly 元（\QFourThreeSavingPlanOnlyPercent\%），紧急购电由\QFourThreePlanOnlyEmergency 降至\QFourThreeEmergencyEnergy\ kWh。

4-3较4-2节省\QFourThreeVsTwo 元（\QFourThreeVsTwoPercent\%），日均节省95\%区间为$[\QFourCIThreeVsTwoLow,\QFourCIThreeVsTwoHigh]$元。经“融合预报、仅计划”这一中间策略分解，0时融合预报与日内调整分别贡献\QFourValueExternalPlan 元和\QFourValueAdjust 元，均以4-2为分母，对应\QFourValueExternalPlanPercent\%和\QFourValueAdjustPercent\%；日内调整占合计节省的\QFourAdjustShareOfSaving\%，其中包含实际状态反馈的作用。

\input{../output/tables/q4_comparison.tex}

\begin{figure}[htbp]
 \centering\includegraphics[width=\textwidth]{../output/figures/q4_costs.pdf}
 \caption{实时价格下的费用、补购与收益来源。a：两族策略的334天现金费用，括号为相邻变化，橙线为增支、蓝线为节省，横轴不从零起；空心点为价格预知启发式对照，不是最优费用界限。b：对应紧急购电量，从零起。c：4-3相对4-2的节省分解，日内调整包含状态反馈。a的误差线对应相邻差，c的误差线对应合计节省，均为配对日均费用差的七日循环块重采样95\%区间乘以334天。}
 \label{fig:q4-costs}
\end{figure}

\FloatBarrier
\subsection{结算效应、预测收益与风险}

图~\ref{fig:q4-costs}采用逐层替换识别费用来源。保持问题二、问题三的决策不变，仅改用实时价结算，分别增支\QFourTwoSettlementEffect 元和\QFourThreeSettlementEffect 元（\QFourTwoSettlementPercent\%与\QFourThreeSettlementPercent\%）；再让策略响应因果价格预测，分别收回\QFourTwoFixedDiff 元和\QFourThreeFixedDiff 元。4-2的日均预测节费95\%区间为$[\QFourCITwoFixedLow,\QFourCITwoFixedHigh]$元。预测利用了部分价格规律，但尚未抵消结算变化带来的全部增支。

更高的预测精度也不必然减少费用或补购。价格预知对照比两族主策略分别少付\QFourTwoPerfectGap 元和\QFourThreePerfectGap 元，但紧急购电均更高，两族分别有105天和109天的日费用高于因果策略。对4-3，取消日内价格修正仅使费用变化\QFourThreeNoIntradayDiff 元，日均差区间$[\QFourCIThreeNoIntradayLow,\QFourCIThreeNoIntradayHigh]$元包含零；信息链修正的费用差同样不显著，且补购略升。这些结果共同表明，点预测的改善只有在改变关键时段的计划购电量与储能配置时，才会形成实际收益。

指定日期结果见附录表~\ref{tab:q42-purchase}至表~\ref{tab:q43-emergency}。两份工作簿\path{result4-2.xlsx}、\path{result4-3.xlsx}分别沿用问题二、问题三的结果口径。全部\QFourStrategyCount 条全年轨迹、规划价格和交付账单通过独立核验，详见\hyperref[sec:model-validation]{第九章}及附录\ref{app:validation-details}。

